\section{评测与运维：把 Agent 变成可管理系统}

\subsection{评测必须覆盖全链路}
Yu Su 提醒：只测最终答案会掩盖大量系统性问题。Agent 评测应覆盖 ``轨迹质量''、``工具调用正确率''、``重规划触发质量'' 和 ``记忆读取准确率''。

\begin{importantbox}{评测口径建议}
\begin{itemize}
    \item Outcome metrics：任务成功率、完成时延、总成本；
    \item Process metrics：中间步骤正确率、无效动作率、反复重试率；
    \item Safety metrics：越权调用率、敏感信息泄露率、错误恢复时间。
\end{itemize}
\end{importantbox}

\begin{knowledgebox}{离线评测与在线评测的分工}
\begin{itemize}
    \item 离线：快速回归算法改动，定位模块问题；
    \item 在线：观测真实分布漂移，校正离线假设；
    \item 两者之间需共享统一事件 schema，否则难以闭环。
\end{itemize}
\end{knowledgebox}

\begin{warningbox}{只看 leaderboard 的风险}
高分可能来自任务模板拟合，而非真实泛化。生产系统必须对分布外任务建立降级策略和人工接管机制。
\end{warningbox}

\subsection{工程可观测性与故障恢复}
稳定 Agent 需要细粒度 tracing：每步决策依据、调用工具、返回结果、记忆命中、计划变更都应可审计。否则定位问题只能靠复现猜测，迭代效率极低。

\subsection{本章小结}
评测和运维不是附加组件，而是 Agent 能否持续进化的基础设施。没有可观测性，就没有可控迭代。

